\documentclass[a4paper,11pt]{article}
\usepackage[T2A]{fontenc}
\usepackage[utf8]{inputenc}
\usepackage[russian]{babel}
\usepackage{amsmath,amssymb,mathtools}
\usepackage{cmbright}
\renewcommand{\familydefault}{\sfdefault}
\usepackage{icomma}
\usepackage[a4paper,left=16mm,right=16mm,top=12mm,bottom=11mm]{geometry}
\usepackage{microtype}
\usepackage{multicol}
\usepackage{xcolor}
\usepackage{hyperref}
\definecolor{accent}{HTML}{15747A}
\definecolor{darkaccent}{HTML}{164C55}
\definecolor{textgray}{HTML}{283238}
\definecolor{muted}{HTML}{65747A}
\definecolor{rulegray}{HTML}{CADADC}
\color{textgray}
\hypersetup{hidelinks}
\setlength{\parindent}{0pt}
\setlength{\parskip}{0pt}
\setlength{\columnsep}{8mm}
\setlength{\abovedisplayskip}{1mm}
\setlength{\belowdisplayskip}{1mm}
\emergencystretch=2em
\newcommand{\sect}[1]{\par\vspace{2.2mm}{\normalsize\bfseries\color{darkaccent}#1}\par\vspace{0.3mm}{\color{rulegray}\rule{\linewidth}{0.45pt}}\par\vspace{0.8mm}}
\newcommand{\fml}[1]{\par\noindent{\normalsize\(\displaystyle #1\)}\par\vspace{2.8mm}}
\begin{document}
\thispagestyle{empty}

\hypersetup{pdftitle={Главные формулы. Анна Трапезникова},pdfauthor={Лёвин Артём Александрович}}
\begin{center}
{\small\bfseries\color{accent}УЧЕБНАЯ ПАМЯТКА \quad\textbullet\quad ОСНОВНЫЕ ФОРМУЛЫ}\par
\vspace{1mm}
{\fontsize{21}{25}\selectfont\bfseries\color{darkaccent} Модуль, неравенства и планиметрия\par}
\vspace{0.8mm}
{\small\color{muted}Анна Трапезникова\quad\textbullet\quad 10 класс · углублённая математика\quad\textbullet\quad 12.10.2026}
\end{center}
\vspace{2mm}
\noindent{\color{accent}\rule{\linewidth}{0.9pt}}
\vspace{1.3mm}
{\bfseries\color{darkaccent}Главные принципы решения}\par
\vspace{1mm}
\begin{multicols}{2}
\raggedcolumns
\small
\noindent\textbullet\quad Модуль раскрывайте по знаку подмодульного выражения.\par\vspace{0.75mm}
\noindent\textbullet\quad Для рациональных неравенств сначала найдите ОДЗ, нули числителя и знаменателя.\par\vspace{0.75mm}
\columnbreak
\noindent\textbullet\quad На числовой прямой учитывайте кратность корней; точки с нулевым знаменателем исключайте.\par\vspace{0.75mm}
\noindent\textbullet\quad В геометрии сначала найдите известный треугольник, затем используйте отношение и площадь.\par\vspace{0.75mm}
\end{multicols}
\vspace{1.1mm}
\noindent{\color{accent}\rule{\linewidth}{0.5pt}}
\vspace{0.6mm}
\begin{multicols}{2}
\raggedcolumns
\normalsize
\sect{Уравнения и неравенства с модулем}
\fml{|A|=a\;(a\ge0)\Leftrightarrow A=\pm a}
\fml{|A|\le a\Leftrightarrow -a\le A\le a}
\fml{|A|\ge a\Leftrightarrow A\le-a\ \lor\ A\ge a}
\sect{Сравнение модулей}
\fml{|A|\le|B|\Leftrightarrow A^2\le B^2}
\fml{|A-B|\le c\;(c\ge0)}
\fml{\Leftrightarrow -c\le A-B\le c}
\columnbreak
\sect{Прямоугольный треугольник}
\fml{a^2+b^2=c^2}
\fml{\sin\alpha=\frac{a}{c},\quad \cos\alpha=\frac{b}{c}}
\fml{\tg\alpha=\frac ab\;(b\ne0)}
\sect{Площадь и вписанный четырёхугольник}
\fml{S_{\triangle}=\frac12 ab\sin\gamma}
\fml{AC\cdot BD=AB\cdot CD+BC\cdot AD}
\end{multicols}
\vfill
{\color{accent}\rule{\linewidth}{0.8pt}}\par\vspace{1.5mm}
{\small\color{muted}Лёвин Артём Александрович эксклюзивно для Анны Трапезниковой}
\end{document}